\section*{第 \ref{ch:g11:becoming-male-female} 章——男性与女性的形成}

\begin{solution}{exo:g11:becoming-male-female:1}
由父亲决定：每个卵子都携带一条 X \omterm{def:g11:cell-cycle-mitosis:chromosome}{染色体}，而一半精子携带 X、一半携带 Y，
因此 XX 与 XY 的可能性相等。
\end{solution}

\begin{solution}{exo:g11:becoming-male-female:2}
两个未分化的性腺，各旁边有两条管道：
一条沃尔夫管（将发育为男性生殖道）和一条米勒管（将发育为女性生殖道）。
两者在 XX 与 XY 胚胎中相同。
\end{solution}

\begin{solution}{exo:g11:becoming-male-female:3}
位于 Y \omterm{def:g11:cell-cycle-mitosis:chromosome}{染色体}上；在约第 7 周时于未分化性腺中表达，
其\omterm{def:g11:gene-expression:protein}{蛋白质}开启一系列\omterm{def:g10:universal-dna:gene}{基因}，
使性腺发育为睾丸。
\end{solution}

\begin{solution}{exo:g11:becoming-male-female:4}
睾酮维持沃尔夫管并使其发育为男性生殖道；
抗米勒管激素使米勒管退化。
\end{solution}

\begin{solution}{exo:g11:becoming-male-female:5}
大脑中的一个中枢释放脉冲式的激素，
使垂体分泌 LH 和 FSH，二者激活性腺。第二性征是指由性激素引起的变化——
体毛、嗓音、乳房、肌肉与脂肪的分布、
生长突增，
以及生殖器官的生长。
\end{solution}

\begin{solution}{exo:g11:becoming-male-female:6}
尽管具有 Y \omterm{def:g11:cell-cycle-mitosis:chromosome}{染色体}，它却发育出了女性生殖道：
男性生殖道需要睾丸分泌的激素，
而女性生殖道在没有任何性腺时也能形成。
\end{solution}

\begin{solution}{exo:g11:becoming-male-female:7}
睾酮使男性管道发育，但不会消除女性管道；
消除女性管道则需要睾丸分泌第二种物质——抗米勒管激素。
\end{solution}

\begin{solution}{exo:g11:becoming-male-female:8}
睾丸形成（Sry 起作用）；沃尔夫管保留，米勒管退化；男性
外生殖器。不育：精子的产生需要 Y \omterm{def:g11:cell-cycle-mitosis:chromosome}{染色体}上的\omterm{def:g10:universal-dna:gene}{基因}，
而转\omterm{def:g10:universal-dna:gene}{基因}未提供这些\omterm{def:g10:universal-dna:gene}{基因}。
\end{solution}

\begin{solution}{exo:g11:becoming-male-female:9}
雌激素：约在 11 岁时达到 50\%，约在 14--15 岁时趋于平稳。睾酮：
约在 13 岁时达到 50\%，约在 16--17 岁时趋于平稳。
\end{solution}

\begin{solution}{exo:g11:becoming-male-female:10}
性激素达到成年水平时，骨骼停止生长；
雌激素比睾酮约早两年达到该水平，
因此女孩的骺板闭合得更早。
\end{solution}

\begin{solution}{exo:g11:becoming-male-female:11}
无睾丸：性腺发育为类似卵巢的性腺；
无睾酮和抗米勒管激素：米勒管发育，沃尔夫管退化，
出生时外生殖器为女性。\omterm{prop:g11:becoming-male-female:puberty}{青春期}时，性腺因缺乏正常的卵子而分泌很少；
若不进行激素治疗，
第二性征发育不良。
\end{solution}

\begin{solution}{exo:g11:becoming-male-female:12}
睾丸形成（\omterm{prop:g11:becoming-male-female:sry}{SRY} 起作用）。沃尔夫管退化（对睾酮无反应）；
米勒管退化（抗米勒管激素起作用）：两种内生殖道均不存在。
外生殖器为女性（其男性形态需要睾酮的作用）。
\omterm{prop:g11:becoming-male-female:puberty}{青春期}时，睾酮转化为雌激素，产生乳房和女性体型，
但无体毛。
\end{solution}

\begin{solution}{exo:g11:becoming-male-female:13}
一旦性腺发育为睾丸或卵巢，其激素便可像在哺乳动物中一样塑造管道和身体，
不同之处在于第一步——
一个对温度敏感的过程取代了 \omterm{prop:g11:becoming-male-female:sry}{SRY}，
成为决定性腺类型的开关。
\end{solution}

\begin{solution}{exo:g11:becoming-male-female:14}
X \omterm{def:g11:cell-cycle-mitosis:chromosome}{染色体}携带数百个每个\omterm{def:g10:cells-common-unit:cell}{细胞}都需要的\omterm{def:g10:universal-dna:gene}{基因}，因此没有胚胎能缺少它而存活；
Y \omterm{def:g11:cell-cycle-mitosis:chromosome}{染色体}除决定睾丸的\omterm{def:g10:universal-dna:gene}{基因}和与精子相关的\omterm{def:g10:universal-dna:gene}{基因}外几乎不携带其他\omterm{def:g10:universal-dna:gene}{基因}，
因此 XX 个体并不缺少任何所需的\omterm{def:g10:universal-dna:gene}{基因}。
Y \omterm{def:g11:cell-cycle-mitosis:chromosome}{染色体}只会传给儿子，
因为女儿从父亲那里得到的是 X \omterm{def:g11:cell-cycle-mitosis:chromosome}{染色体}。
\end{solution}

\begin{solution}{exo:g11:becoming-male-female:15}
这意味着：当没有睾丸激素参与时，女性生殖道便会发育——
这说明的是需要哪些信号，而非高低之分。
它反映了一个共同的胚胎蓝图，其中一个\omterm{def:g10:cells-common-unit:cell}{细胞}谱系即睾丸，
借助两种激素施加另一种发育方向：
两性是同一程序的两种结果，二者只相差一个开关。
\end{solution}

\begin{solution}{pb:g11:becoming-male-female:1}
\textbf{1.} 卵巢对女性生殖道并非必需，女性生殖道在没有任何性腺时也能形成；
而睾丸对男性生殖道的形成
以及女性管道的消除是必需的。

\textbf{2.} 第 22 天后，管道已经开始分化，
且已受性腺影响；实验必须
在这一决定作出之前施加作用。

\textbf{3.} 睾丸的激素经血液运送至管道：
移植物远离管道，但仍能使其男性化。

\textbf{4.} 睾酮促使男性管道发育；它并不消除
女性管道。

\textbf{5.} 另一种睾丸物质必须使苗勒管
退化；可从睾丸中提取该物质并单独植入以作验证，
预期女性管道消失，男性管道
不形成——这正是 AMH 的作用。

\textbf{6.} 手术一侧既无睾酮也无 AMH：仅有女性管道，
无男性管道。未手术一侧：有男性管道，女性管道已退化。

\textbf{7.} 移植一侧：有男性管道，无女性管道；
另一侧：仅有女性管道。

\textbf{8.} 晶体一侧：两种管道均存在；另一侧
仅有女性管道。

\textbf{9.} 每个胎儿都可作为自身对照：两侧
之间的差异只能源于睾丸的局部存在、
睾酮的存在或缺失。

\textbf{10.} 其卵巢被摘除或被压制，而
来自另一动物的移植睾丸在异体宿主中不产生精子；
生殖道为雄性，但不产生配子。

\textbf{11.} 存在 \omterm{prop:g11:becoming-male-female:sry}{SRY}：形成睾丸；睾酮与 AMH：男性生殖道、
男性解剖结构；\omterm{prop:g11:becoming-male-female:puberty}{青春期}睾酮：男性特征。不育
（缺少产生精子所需的 Y \omterm{def:g11:cell-cycle-mitosis:chromosome}{染色体}\omterm{def:g10:universal-dna:gene}{基因}）。

\textbf{12.} \omterm{prop:g11:becoming-male-female:sry}{SRY} 无功能：无睾丸；无激素：女性生殖道与
解剖结构；性腺发育不良。

\textbf{13.} 有睾丸；沃尔夫管发育（睾酮起作用）；
苗勒管也保留（对 AMH 无反应）：男性个体，
另有子宫和输卵管；男性解剖结构与\omterm{prop:g11:becoming-male-female:puberty}{青春期}。

\textbf{14.} 只有第三种：睾丸与男性管道功能正常，
尽管保留的子宫可能阻碍睾丸下降。
第一种缺少产生精子的\omterm{def:g10:universal-dna:gene}{基因}；第二种性腺无功能。

\textbf{15.} \omterm{def:g11:cell-cycle-mitosis:chromosome}{染色体}性别只有经由中间步骤
才能预测解剖性别：\omterm{prop:g11:becoming-male-female:sry}{SRY} 功能正常、激素与受体功能正常。
改变任一环节，两个层次即会分离。

\textbf{16.} 女孩：10 岁时达 20\%，约 11 岁开始生长突增，约 15 岁停止。
男孩：约 12.5 岁时达 20\%，约 13.5 岁开始生长突增，约 17 岁停止。

\textbf{17.} 女孩约 4 年，男孩约 3.5 年——但男孩
起点较高，因为童年期生长速率持续
多两年。

\textbf{18.} 男孩约 $3.5 \times 8 = \qty{28}{cm}$，女孩 $4 \times
7 = \qty{28}{cm}$：两者突增幅度相似；成年身高
差异主要来自男孩突增前多出的
两年童年期生长。

\textbf{19.} 过早升高会使骺板
提前数年闭合：孩子现在长得快，但成年身高矮小。
延迟激素作用则延长童年期生长。

\textbf{20.} 睾酮促使男性管道发育，
抗苗勒管激素（AMH）则消除女性管道；二者均由
睾丸产生，睾丸这一器官由单一\omterm{def:g10:universal-dna:gene}{基因} \omterm{prop:g11:becoming-male-female:sry}{SRY} 开启。
\end{solution}
